\clearpage
\SourceLeaf{207}
\renewcommand{\BookRunningTitle}{Handschriftlicher Nachlass}
\thispagestyle{empty}
\pdfbookmark[0]{Handschriftlicher Nachlass}{handschriftlicher-nachlass}
\vspace*{.25\textheight}
\begin{center}
{\Large HANDSCHRIFTLICHER}\par\vspace{2.8cm}
{\Huge NACHLASS.}
\end{center}
\vfill

\clearpage
\SourceLeaf{209}
\renewcommand{\BookRunningTitle}{Nachlass. Lehre von den Resten}
\pdfbookmark[0]{Lösung einer reinen Kongruenz. Lehre von den Resten}{analysis-residuorum-pars-prior}
\begin{center}
{\large LÖSUNG DER KONGRUENZ $X^m-1\equiv0$.}\par\medskip
LEHRE VON DEN RESTEN. SECHSTES KAPITEL. ERSTER TEIL.
\end{center}
\BookRule{.18\linewidth}
\renewcommand{\Article}[1]{\par\medskip\Needspace{4\baselineskip}\pdfbookmark[1]{Art. #1}{analysis-residuorum-art-#1}\centerline{\textbf{#1.}}\nobreak\smallskip}
\Article{237}
In Kapitel \textsc{iii} haben wir gezeigt, dass die Kongruenz $x^n\equiv1$, wenn eine Primzahl $p$ als Modul genommen wird, $\mu$ Wurzeln hat, wobei $\mu$ der größte gemeinsame Teiler der Zahlen $n$ und $p-1$ ist, und dass diese Wurzeln mit denen der Kongruenz $x^\mu\equiv1$ genau übereinstimmen. Es genügt daher, den Fall zu betrachten, dass $n$ ein Teiler von $p-1$ ist. Dass sich aber nicht nur die Lösung der Kongruenz $x^n-1\equiv0$, sondern auch die jeder anderen Kongruenz für beliebige Moduln aus ihrer Lösung für Primzahlmoduln ableiten lässt, wurde schon an verschiedenen Stellen gezeigt und wird weiter unten (Kapitel \textsc{viii}) ausführlicher entwickelt.

\Article{238}
Doch auch hier brauchen wir nicht stehen zu bleiben. In demselben Kapitel \textsc{iii} haben wir dargelegt, dass die Lösung der Kongruenz $x^n\equiv1$ von der Lösung ähnlicher Kongruenzen $x^a\equiv1$, $x^b\equiv1$ usw. abhängt, wobei $a,b$ usw. Primzahlen oder Primzahlpotenzen sind und $n$ das Produkt dieser Zahlen ist. Sind nämlich $A,B$ usw. jeweils beliebige Wurzeln der Kongruenzen $x^a\equiv1$, $x^b\equiv1$ usw., so ist ihr Produkt $AB\ldots$ eine der Wurzeln der Kongruenz $x^n\equiv1$. Unsere Untersuchungen beschränken sich daher auf die Lösung der Kongruenz $x^n\equiv1\;(\mathrm{mod.}\,p)$, wobei $p$ eine Primzahl und $n$ eine Primzahl oder Primzahlpotenz und zugleich ein Teiler von $p-1$ ist.

\SourceLeaf{210}
\Article{239}
Aus Kapitel \textsc{iii} ist ferner bekannt, dass es unter den Wurzeln der Kongruenz $x^n\equiv1$ immer solche gibt, durch deren Potenzen sich alle übrigen darstellen lassen. Bezeichnet also $r$ eine solche Wurzel (wir haben sie oben \emph{primitiv} genannt, wenn $n=p-1$, und behalten diesen Ausdruck hier bei, wenn auch in weiterem Sinn), so sind alle Wurzeln der gegebenen Kongruenz
\[
1,\ r,\ rr,\ r^3\quad\dotsfive\quad r^{n-1}
\]
Wir müssen daher alle Anstrengungen darauf richten, solche Wurzeln zu finden, da sich die übrigen von selbst ergeben, sobald diese gefunden sind. Der Kürze halber bezeichnen wir jede Potenz von $r$ durch ihren in Klammern eingeschlossenen Exponenten. Dabei bezeichnet $(0)$ die Eins, $(1)$ eine beliebige primitive Wurzel der Kongruenz $x^n\equiv1$, $(2)$ das Quadrat von $(1)$ usw., sodass die Folge $(0),\ (1),\ (2),\ (3),\dotsfour(n-1)$ alle Wurzeln umfasst. Ferner ist bekannt, dass $(k)$ immer eine solche primitive Wurzel ist, wenn $k$ zu $n$ teilerfremd ist, also in unserem Fall, in dem $n$ eine Potenz der Primzahl $t$, $=t^\nu$, ist, immer dann, wenn $t$ die Zahl $k$ nicht teilt. Die Zeichen $(1)$, $(2)$ usw. sind für sich genommen offensichtlich unbestimmt. Sobald aber $(1)$ ein bestimmter Wert gegeben wird, sind auch alle übrigen bestimmt.

\Article{240}
Da wir vor den übrigen vor allem die primitiven Wurzeln untersuchen wollen, müssen wir diese zunächst von den anderen trennen. Das geschieht, indem wir aus der Folge $(0),(1),(2)\ldots(n-1)$ alle Terme $(k)$ entfernen, bei denen $k$ durch $t$ teilbar ist. Ist aber $n$ eine Primzahl oder $\nu=1$, so ist nur $(0)$ zu entfernen. Bevor wir die verbleibenden Wurzeln untersuchen, empfehlen wir dem Leser nachdrücklich, sich einige Beispiele zu bilden, damit er alles, was ohne diese vielleicht zu allgemein gehalten erschiene, konkret vor sich sieht. Wir fügen ein Beispiel hinzu; dadurch wird es aber nicht überflüssig, andere selbstständig auszuarbeiten.

Es sei $p=29$; $n=7$. Die sieben Wurzeln der Kongruenz $x^7\equiv1\;(\mathrm{mod.}\,29)$ sind dann $1,7,16,20,23,24,25$. Da $n$ eine Primzahl ist, sind alle diese Wurzeln außer $1$ primitiv. Setzt man also $7=(1)$, so haben die Zeichen folgende Bedeutungen:
\[
\begin{array}{ccccccc}
(0)&(1)&(2)&(3)&(4)&(5)&(6)\\
1&7&20&24&23&16&25
\end{array}
\]

\SourceLeaf{211}
Man wird sich außerdem erinnern, dass die Zeichen $(n)$ und $(0)$, $(n+1)$ und $(1)$ usw. und allgemein $(a)$ und $(b)$ gleichbedeutend sind, wenn $a\equiv b\;(\mathrm{mod.}\,n)$.

\Article{241}
Für unser Vorhaben müssen wir aber noch auf eine andere Weise vorgehen. Wir behalten nämlich nur die Terme $(k)$, bei denen $k$ nicht durch $t$ teilbar ist; ihre Anzahl ist $\frac{t-1}{t}\pdot n=\lambda$. Alle diese Zahlen oder ihnen modulo $n$ kongruente Zahlen lassen sich durch aufeinanderfolgende Potenzen einer Zahl darstellen. Diese sei $=\rho$; dann werden alle primitiven Wurzeln der Kongruenz $x^n\equiv1$ so bezeichnet:
\[
(1)\quad(\rho)\quad(\rho^2)\quad(\rho^3)\ \dotsfour\ (\rho^{\lambda-1})
\]
Durch dieses Verfahren erreichen wir, dass sämtliche nicht primitiven Wurzeln vollständig ausgeschlossen werden. Die Gründe und Vorteile werden weiter unten deutlicher. In unserem Beispiel können wir also $\rho=3$ setzen; die primitiven Wurzeln der Kongruenz $x^7\equiv1$ sind dann folgendermaßen angeordnet:
\[
\begin{array}{rcccccc}
&(1)&(3)&(3^2)&(3^3)&(3^4)&(3^5)\\
\text{oder}&(1)&(3)&(2)&(6)&(4)&(5)\\
\text{das sind}&7&24&20&25&23&16
\end{array}
\]

\Article{242}
Damit der Leser das Ziel der folgenden Untersuchungen kennt, ist es zweckmäßig, den Satz anzugeben, den wir beweisen und erläutern wollen.

\emph{Hat die Zahl $\lambda$, die $=t^{\nu-1}\pdot t-1$ ist, die Primfaktoren $a,b,c,d$ usw., und ist $\lambda=a^\alpha b^\beta c^\gamma\ldots$, so hängt die Lösung der Kongruenz $x^n-1\equiv0$ von der Lösung von $\alpha+\beta+\dotstwo$ Kongruenzen niedrigeren Grades ab, von denen $\alpha$ den Grad $a$, $\beta$ den Grad $b$, $\gamma$ den Grad $c$ haben usw.}

In unserem Beispiel hängt also die Lösung der Kongruenz $x^7\equiv1$ von einer Kongruenz zweiten und einer weiteren dritten Grades ab. Allgemein erkennt man, dass der Grad dieser Kongruenzen niemals vom Modul $p$ abhängt. Um zum Beweis dieses Satzes zu gelangen, müssen wir aber einige Sätze über den Zusammenhang zwischen Kongruenzen und ihren Wurzeln vorausschicken, obwohl diese Untersuchungen eigentlich erst im achten Kapitel weiterzuverfolgen sind.

\SourceLeaf{212}
\Article{243}
\textsc{Satz.}\quad\emph{Die Kongruenz}
\[
x^m+Ax^{m-1}+Bx^{m-2}+\dotsfour+N\equiv0\;(\text{Primzahlmodul})
\]
\emph{sei so beschaffen, dass nach dem Ausmultiplizieren des Produkts aus den $m$ Faktoren $x-r$, $x-r'$, $x-r''$, $x-r'''\ldots$, das $x^m+ax^{m-1}+bx^{m-2}\ldots+n$ ergebe, die Beziehungen $A\equiv a$, $B\equiv b$, $C\equiv c$ usw. modulo $p$ gelten. Dann sind die Größen $r,r',r''\ldots$ Wurzeln der gegebenen Kongruenz, und sie hat keine anderen.}

\emph{Beweis.}\quad I. Es gilt immer
\[
\begin{aligned}
&x^m+Ax^{m-1}+Bx^{m-2}+\ldots\\
&\qquad\equiv x^m+ax^{m-1}+bx^{m-2}+\ldots\;(\mathrm{mod.}\,p)
\end{aligned}
\]
Die rechte Seite der Kongruenz wird aber $=0$, wenn man $x=r$, $x=r'$, $x=r''$ usw. setzt. Für diese Werte von $x$ wird daher die linke Seite $\equiv0\;(\mathrm{mod.}\,p)$. Damit ist der erste Teil bewiesen.

II. Würde noch ein weiterer Wert $\rho$, der keiner der Zahlen $r,r'$ usw. kongruent ist, die gegebene Kongruenz erfüllen, so wäre
\[
\begin{aligned}
0&\equiv\rho^m+A\rho^{m-1}+B\rho^{m-2}+\ldots\\
&\equiv\rho^m+a\rho^{m-1}+b\rho^{m-2}+\ldots\\
&\equiv(\rho-r)(\rho-r')(\rho-r'')(\rho-r''')\ldots
\end{aligned}
\]
Da aber keiner der Faktoren $\rho-r$, $\rho-r'$, $\rho-r''$ usw. $\equiv0$ ist, ist es unmöglich, dass ihr Produkt $\equiv0$ ist, weil $p$ eine Primzahl ist. Außer den Wurzeln $r,r'$ usw. gibt es also keine anderen. Damit ist der zweite Teil bewiesen.

\Article{244}
\textsc{Aufgabe.}\quad\emph{Es seien $r,r',r''\ldots$ unbekannte Größen, deren Anzahl $=m$, deren Summe $=\alpha$, deren Summe der Quadrate $=\beta$, deren Summe der Kuben $=\gamma\ldots$ und deren Summe der Potenzen mit Exponenten $m$ gleich $=\mu$ sei. Gegeben seien aber nicht diese Zahlen selbst, deren Anzahl ebenfalls $=m$ ist, sondern andere Zahlen $\alpha',\beta',\gamma'$ usw., die ihnen einzeln modulo $p$ kongruent sind, wobei dieser Modul eine Primzahl und $>m$ sei. Gesucht ist eine Kongruenz des Grades $m$, deren Wurzeln $r,r',r''$ usw. sind.}

\emph{Lösung.}\quad Man betrachte $r,r',r''$ usw. als Wurzeln einer Gleichung
\[
x^m+Ax^{m-1}+Bx^{m-2}+Cx^{m-3}+\ldots=0
\]
und bestimme ihre Koeffizienten $A,B,C$ usw. nach der bekannten Methode, wobei nur die Kongruenz an die Stelle der Gleichheit tritt. Man setze also
\SourceLeaf{213}
\[
\begin{aligned}
-A&\equiv\alpha'\\
-2B&\equiv\beta'+A\alpha'\\
-3C&\equiv\gamma'+A\beta'+B\alpha'\\
-4D&\equiv\delta'+A\gamma'+B\beta'+C\alpha'\\
&\text{usw.}\\
-mN&\equiv\mu'+A\lambda'+\text{usw.}
\end{aligned}
\]
\begingroup\rightskip=0pt plus2em\relax Diese Koeffizienten können nicht unbestimmt sein, weil alle Zahlen \mbox{$1,2,3\ldots m<p$}. Ich behaupte, dass die Kongruenz
\[
x^m+Ax^{m-1}+Bx^{m-2}+\ldots+N\equiv0
\]
die gesuchte ist.\par\endgroup

\emph{Beweis.}\quad Die Gleichung, deren Wurzeln $r,r',r'',r'''$ usw. sind, sei
\[
x^m+ax^{m-1}+bx^{m-2}+\ldots=0
\]
Dann gilt
\[
\begin{aligned}
-a&=\alpha\\
-2b&=\beta+a\alpha\\
-3c&=\gamma+a\beta+b\alpha\\
-4d&=\delta+a\gamma+b\beta+c\alpha\\
&\text{usw.}
\end{aligned}
\]
\Needspace{12\baselineskip}Daraus wird jedem klar, dass
\[
a\equiv A,\quad b\equiv B,\quad c\equiv C\ \text{usw.}\;(\mathrm{mod.}\,p)
\]
Also sind nach dem vorigen Abschnitt die Zahlen $r,r',r''$ usw., die Wurzeln der \emph{Gleichung}
\[
x^m+ax^{m-1}+bx^{m-2}+\ldots=0
\]
sind, zugleich Wurzeln der \emph{Kongruenz}
\[
x^m+Ax^{m-1}+Bx^{m-2}+\ldots\equiv0\qquad\text{Q. E. D.}
\]
Wir überlassen es den Lesern, sich Beispiele zu bilden.

\Article{245}
Wir kehren zu unserem Vorhaben zurück. Mit den in den Abschnitten 242 und zuvor verwendeten Zeichen wollen wir zeigen: Ist $\lambda$ das Produkt beliebiger Faktoren
\SourceLeaf{214}
$efg$ usw., so können die primitiven Wurzeln der Kongruenz $x^n\equiv1$, deren Anzahl $\lambda$ ist, so in $e$ Klassen eingeteilt werden, dass die Summen der jeweils derselben Klasse zugeordneten Wurzeln durch eine Kongruenz des Grades $e$ gegeben werden. Setzt man diese als bekannt voraus, so lässt sich jede Klasse weiter so in $f$ Gruppen unterteilen, dass die Summe jeder Gruppe durch eine Kongruenz des Grades $f$ gegeben wird. Diese Gruppen lassen sich wiederum unterteilen usw., bis man zu den einzelnen Wurzeln gelangt.

\Article{246}
\emph{Definition.}\quad Die Gesamtheit \emph{aller} Terme, die in einer solchen Form $(\rho^{ke+\alpha})$ (Abschnitt 241) enthalten sind, nennen wir eine \emph{vollständige Periode} oder einfach eine \emph{Periode}. Dabei bezeichnet $e$ einen beliebigen Teiler von $\lambda$, $\alpha$ eine beliebige gegebene Zahl und $k$ alle ganzen Zahlen von $0$ bis $\frac{\lambda}{e}-1$. Der Kürze halber bezeichnen wir eine solche Periode mit $(e\ast\alpha)$. So bilden in unserem Beispiel die Terme
\[
\begin{array}{rclcl}
&(1),(2),(4)&\text{Perioden}&(2\ast0)&\text{,}\\
&(3),(6),(5)&&(2\ast1)&\\
\text{diese aber}&(1),(6)&\text{diese}&(3\ast0)&\\
&(3),(4)&&(3\ast1)&\\
&(2),(5)&&(3\ast2)&
\end{array}
\]

\emph{Wenn nun alle Terme auf beliebige Weise in Perioden verteilt werden, die einzelnen Perioden wiederum in kleinere Perioden usw., so behaupten wir, dass damit das im vorigen Abschnitt Versprochene erreicht wird.}

Bevor wir diese Darstellung selbst beginnen, zeigen wir jedoch, dass die Bildung einer solchen Periode nichts Unbestimmtes enthält, obwohl sie von den beiden in gewisser Weise willkürlichen Größen $r,\rho$ abhängt. Wie auch immer diese Größen gewählt werden, kommen stets dieselben Terme in dieselbe Periode, sofern vorgeschrieben ist, wie viele Terme die Periode enthalten soll.

Das Kriterium dafür, dass zwei Terme $A,B$ in derselben Periode liegen, ergibt sich daraus, dass beide in einer solchen Form $(\rho^{ke+\alpha})$ enthalten sind, also $A\equiv r^{\rho^{ke+\alpha}}$, $B\equiv r^{\rho^{k'e+\alpha}}\;(\mathrm{mod.}\,p)$. Dabei ist $r$ eine primitive Wurzel der Kongruenz $x^n\equiv1\;(\mathrm{mod.}\,p)$ und $\rho$ eine primitive Wurzel der Kongruenz $x^\lambda\equiv1\;(\mathrm{mod.}\,n)$; siehe oben.

Zu beweisen ist: Werden anstelle der Zahlen $r,\rho$ andere, etwa $s,\sigma$, gewählt, so sind $A$ und $B$ in ähnlichen Formen $s^{\sigma^{le+\beta}}$, $s^{\sigma^{l'e+\beta}}$ enthalten.

\SourceLeaf{215}
Es sei $s^m\equiv r\;(\mathrm{mod.}\,p)$; $\sigma^\mu\equiv\rho\;(\mathrm{mod.}\,n)$ und $m\equiv\sigma^\zeta\;(\mathrm{mod.}\,n)$, was möglich ist, weil $r,\rho$ primitive Wurzeln sind. Dabei ist $m$ zu $n$ und $\mu$ zu $\lambda$ teilerfremd (Kapitel \textsc{iii}). Durch die entsprechenden Substitutionen erhalten wir
\[
A\equiv s^{\sigma^{\mu ke+\mu\alpha+\zeta}},\qquad B=s^{\sigma^{\mu k'e+\mu\alpha+\zeta}}\qquad\text{Q. E. D.}
\]
